%  07_testing.tex  –  Chapter 6: Testing


\chapter{Testing}
\label{ch:testing}

\section{Unit Testing}
\label{sec:unit-testing}

\subsection{Hardware Components}
\label{subsec:hw-testing}

\subsubsection*{DHT11 Temperature and Humidity Sensor}

Over the recorded experiment sessions, the DHT11 was reliable in its operations as far as the
physical connections were concerned. There were a few sessions that showed values of
0$^\circ$C and 0\,\% humidity. This was due to a loose cable connection on the breadboard
that was corrected after the hardware experiment session. It can also be caused by the sensor
picking up the heat from the other ones, but it was moved further away to avoid this scenario.
On the June 12th session for example, temperature stayed within a 1.8$^\circ$C range
(27.3--25.5$^\circ$C) and humidity within 6\,\% (39--45\,\%) across a window of almost
24 hours, indicating stable and credible readings. Only once on this day, the sensor read 0
values, but it was solved without external interaction, through the retry mechanism.

The DHT11 retry mechanism was triggered only once during the valid sessions (not the test
sessions), which confirms that when the physical connection is secure, the sensor returns valid
data on the first attempt.

\subsubsection*{MQ Gas Sensors and ADC}

The MQ sensors were evaluated across multiple valid measurement sessions conducted in the
Aradului neighbourhood of Timișoara. The readings were collected under similar conditions: an
indoor environment with a nearby open window (most of the time), at temperatures of
24$^\circ$C to 30$^\circ$C.

\textit{Pre-calibration and post-calibration comparison}

On June 6 the sensors were fully calibrated. In order to test if it was correct, readings
collected on May 25th (before) have been compared with those taken on June 13 (after). Both
sessions were taken indoors, in the same space. Only difference is that on May 25th the window
was open, while on June 13th wasn't. Also the first recordings were taken at 26.6--26.9$^\circ$C
and 48\,\% humidity, while the latter at 25.1--26.4$^\circ$C and 40--48\,\% humidity.

\vspace{12pt}
\begin{table}[H]
  \centering
  \caption{Table 5: PPM values (before recalibration and after)}
  \label{tab:calibration-compare}
  \renewcommand{\arraystretch}{1.3}
  \begin{tabular}{lrrr}
    \toprule
    \textbf{Sensor} & \textbf{May 25\textsuperscript{th} PPM mean} & \textbf{June 13\textsuperscript{th} PPM mean} & \textbf{Change} \\
    \midrule
    MQ135 &    267.81 & 1{,}081.04 & $+304\,\%$ \\
    MQ3   &      0.58 &      0.28  & $-52\,\%$  \\
    MQ6   & 1{,}299.27 & 1{,}184.39 & $-9\,\%$  \\
    MQ7   &    150.46 &     98.95  & $-34\,\%$  \\
    MQ8   & 1{,}152.35 & 1{,}354.16 & $+17\,\%$ \\
    \bottomrule
  \end{tabular}
\end{table}
\vspace{12pt}

Even if the values from the MQ7 decreased quite a bit, they are still above the threshold.
Values for MQ3 (alcohol/benzene) decreased as expected due to the lower sensitivity of $R_O$
in recalibration. Values for MQ135 increased quite a bit on June 13th. This is because during
that session, the window was closed the whole time, while people were in the room, so the
CO\textsubscript{2} accumulated. The readings are accurate, since the new $R_O$ value results
in a change in the formula ($R_S/R_O$), which leads to lower PPM.

The June 13th MQ7 readings of around 99\,ppm are higher than the warning level of 35\,ppm. It
is explained in Section 6.3 why this warning limit is exceeded, even after calibration. The
accuracy of MQ sensors is somewhere around $\pm$15--30\,\%. Also the MQ7 sensor has
difficulties in detecting values which are so low.

\textit{June 13 reading stability}

The June 13th session goes for approximately 6.5 hours. Below is shown the stability of
readings across the session by comparing the first and last readings (03:12 to 09:44):

\vspace{12pt}
\begin{table}[H]
  \centering
  \caption{Table 6: June 13\textsuperscript{th} stability comparison}
  \label{tab:stability}
  \renewcommand{\arraystretch}{1.3}
  \begin{tabular}{lrrr}
    \toprule
    \textbf{Sensor} & \textbf{03:12 reading (ppm)} & \textbf{09:44 reading (ppm)} & \textbf{Variation} \\
    \midrule
    MQ3 & 1{,}022.61 & 1{,}133.97 & $+10.9\,\%$ \\
    MQ6 & 1{,}094.34 & 1{,}238.04 & $+13.1\,\%$ \\
    MQ7 &     93.80  &    101.98  & $+8.7\,\%$  \\
    MQ8 & 1{,}330.29 & 1{,}375.58 & $+3.4\,\%$  \\
    \bottomrule
  \end{tabular}
\end{table}
\vspace{12pt}

The increase in readings indicates slow buildup of gases over time in an indoor space with no
ventilation. The sensor which displays the lowest change is MQ8, while the one with the
highest change is MQ6. Nothing unusual since the closed space will, in time, gather a higher
concentration of gases.

\textit{Temperature and humidity correction factor}

It's been confirmed that the correction factor used for $R_S$ before PPM calculations was
providing results. During the May 25th session, when the temperature was 26.6$^\circ$C and
humidity about 48\,\%, the value of the correction factor is around 0.977. For the June 13th
session, when the temperature is around 26$^\circ$C and humidity is 47\,\%, correction is
around 0.978.

\textit{ADC saturation}

The MQ6 sensor reached the maximum ADC value of 4095, corresponding to 3.3\,V, in the earlier
unstructured test sessions captured in the May 2026 data with no DHT11 readings. While those
sessions are excluded from quantitative analysis, the saturation event confirms that the
current load resistor configuration limits the measurable range of MQ6 at high concentrations.
In the valid test sessions, MQ6 reached a maximum raw value of 1{,}515 (May 24), 1{,}467
(May 25), and 1{,}443 (May 28), all well below saturation, indicating that the sensor behaves
within its measurable range under the observed indoor conditions.

\textit{Extreme event observation}

One important incident took place on June 5 when the following levels were obtained: MQ135
attained a level of 31{,}019\,ppm, MQ6 got a value of 36{,}848\,ppm, MQ7 acquired a level of
3{,}093\,ppm and MQ8 acquired 2{,}939\,ppm. This particular experiment is not included in
quantitative analyses because of the accompanying fault in DHT11. However, the values obtained
demonstrate the effectiveness of the system for recognizing an increase in gas concentration
and activating the alarm system properly. It should also be noted that this session predates
the calibration procedure mentioned further in this section, and the initial $R_O$ values were
used, which turned out to be too high due to inadequate preheating of sensors.

\textit{Sensor Recalibration and Validation}

$R_O$ values employed in the calculation of the initial PPM (Table 2, initial values) have
been derived from a baseline measurement done following approximately 14 hours of uninterrupted
sensor work. Preheating for 24 to 48 hours is recommended in the MQ series sensor manual prior
to calibration for the first time for achieving the equilibrium between the heating element and
sensing element inside the sensor. Consequently, since the initial baseline value has been
acquired prior to the attainment of this balance, resulting $R_O$ values were inflated causing
the PPM values to be unreasonably higher, such as readings for MQ7 sensors at around 400 to
600\,ppm, and sometimes even for MQ135 sensors at more than 10{,}000\,ppm.

\textit{Recalibration Process}

Following a second day of continuous use, making a total of four days in line with the
recommended datasheet value for preheating time, a new baseline reading was performed using the
``Aradului ap.'' station in Timișoara over the course of about 10 hours using 19 readings. The
final two readings were dropped in this case due to malfunctioning of the DHT11 sensor, which
exhibited its failure mode in the form of 0.0$^\circ$C and 0\,\% humidity readings while
caching data from previous readings.

The $R_O$ values for all sensors were recalculated based on the average voltage readings using
the same corrections for the voltage divider and resistance as before, represented in Eq.\ (3)
and Eq.\ (2).

\vspace{12pt}
\begin{table}[H]
  \centering
  \caption{Table 7 -- $R_O$ recalibration comparison}
  \label{tab:ro-recalibration}
  \renewcommand{\arraystretch}{1.3}
  \begin{tabular}{lrrr}
    \toprule
    \textbf{Sensors} & \textbf{Old $R_O$ ($\Omega$)} & \textbf{New $R_O$ ($\Omega$)} & \textbf{Change} \\
    \midrule
    MQ135 & 47{,}059 & 30{,}688 & $-35\,\%$ \\
    MQ3   &  9{,}870 &  7{,}524 & $-24\,\%$ \\
    MQ6   &  8{,}880 &  8{,}490 & $-4\,\%$  \\
    MQ7   & 28{,}712 & 11{,}539 & $-60\,\%$ \\
    MQ8   & 54{,}886 & 38{,}502 & $-30\,\%$ \\
    \bottomrule
  \end{tabular}
\end{table}
\vspace{12pt}

In addition to the previously taken measurements with the window open, another session was
taken with the window closed using the recalculated $R_O$ values. The readings for the MQ7
sensor dropped from the 400--600\,ppm range to only 51--76\,ppm, an important drop, although
still higher than 35\,ppm. In contrast, the voltage of the MQ135 sensor (which is eCO2)
increased by around 28\,\% when the window was closed, which is expected from the increase in
CO\textsubscript{2} levels from breathing in a closed room. Voltage dropped by around 24\,\%
in the MQ3 sensor (which measures alcohol and benzene), while the other readings remained
relatively stable.

\textit{MQ7 Persistent Threshold Exceedance}

After recalibration, MQ7 still consistently reads 51--76\,ppm, above the 35\,ppm threshold,
even with no obvious CO source present.

A few factors likely explain this. MQ-series sensors have a manufacturer-stated accuracy
tolerance of around $\pm$15--30\,\% even when properly calibrated, so a true reading near the
threshold can easily appear 50\,\% higher. This is because the 35\,ppm limit is also determined
by the OSHA 8 hours workplace exposure limit, which is a conservative industrial standard and
not a danger level for a residential building. Normal levels of indoor CO, when there is no
source of combustion, are significantly lower, but MQ7 is less accurate because it is intended
to detect higher concentrations. Finally, the recalibration of MQ7 was made according to only
one session of 19 measurements, so there could be still some offset in MQ7 zero-level reading,
which will result in all values being shifted upwards.

Taking into account all above mentioned facts, it was decided to leave the warning signal
unchanged and not to change the warning threshold even more. The warning means that the system
correctly detects that MQ7 readings are increased in comparison with the clean air baseline
value, which is its correct task.

\subsubsection*{Sound Sensor}

The KY-037 analog output was reliable for all recordings. Values of \texttt{Sound\_raw} ranged
from 111 to 117 during May recordings, equivalent to voltages between 0.089 and 0.094\,volts.
DOUT interrupt sound event counter showed 0 during all the recording segments taken. This was
expected, considering that the place where the station was running at the time was quiet.

\subsubsection*{Warning System and LED Indicators}

The warning mechanism test was done indirectly for all the three gas events that were valid.
For every event, the high PPM values triggered a warning notification that identified the
sensors that had readings which went above the threshold set. The warning message was captured
in the CSV file and in the data that went to the server. The red LED was observed to light for
5 seconds at the end of each cycle where a warning was generated, then turned off. The blue LED
turned on whenever a reading was sent, to show that the data was sent.

\subsection{Software Components}
\label{subsec:sw-testing}

\subsubsection*{PPM Computation}

Testing of calculation of PPM through the pipeline takes place with help of recalibrated
baselines presented in section 6.1.1. $R_S$ for each sensor that has been received according
to the formula for calculating voltage~(3) and resistance formula~(2) in terms of the
recalibration baseline with regard to the average value of voltage corresponds to $R_O$ value
for each sensor. This guarantees right usage of the voltage divider correction formula and
$R_S$ value formula in the script. Also, it is guaranteed that the recalibrated $R_O$ value is
consistent with the formula used for calculation when the sensor operates.

\subsubsection*{Local CSV Logging}

Function of local CSV logging was tested with the restriction of the connection with the remote
server where CSV files usually should be uploaded to. Regardless of that, it operates without
any problems, completes measuring cycle and adds new data line into the local CSV file without
any failures. HTTP POST operation failure is processed by exception handler, while script
itself doesn't stop working.

\subsubsection*{Station Toggle}

The button-based toggle was tested by pressing the button during an active data collection
session. The \texttt{run\_event} flag was correctly cleared and data collection paused.
Pressing the button again resumed collection.

\subsubsection*{External API Integration}

Two external APIs are called during each data collection cycle: \texttt{ipinfo.io} for device
location and OpenWeatherMap for weather symbols~\cite{openweathermap}. They were verified to
function correctly across all recorded sessions. \texttt{ipinfo.io} API identified the
station's location as Timișoara, Romania, which was the right one. OpenWeatherMap API returned
``Sunny'' for all May sessions (since all of those sessions were sunny) and ``Rainy'' for the
June sessions (only the ones which showed signs of rain or straight rain).

\section{Integration Testing}
\label{sec:integration-testing}

This way it was verified the data pipeline. From sensor readings until what is displayed on
the website and mobile app.

\vspace{12pt}
\begin{figure}[H]
  \centering
  \includegraphics[width=0.65\textwidth]{integration_flow}
  \caption{Figure 41: Integration test data flow}
  \label{fig:integration-flow}
\end{figure}
\vspace{12pt}

\subsubsection*{End-to-End Data Flow}

Once the station was operational and connected to the network, it was confirmed that the
readings of the sensors are entered into the PostgreSQL database and can be seen on the website
in a matter of seconds. Each time a cycle of 30 minutes is completed, the Raspberry Pi makes a
POST request to \texttt{/api/sensor}. Then the Flask server enters the row right away, once it
gets the request.

\subsubsection*{Website Authentication}

The authentication flow process has been tested by using the following methods. Registration
through a valid Gmail account followed by activation through the link received in the inbox.
Logging in before account activation resulted in the appropriate error message. Login with
wrong credentials resulted in the error message. For the forgot password flow, the reset link
was successfully sent. Upon resetting the password, login could be done with the new password.
Dummy user accounts have been registered in the database and validated, proving the entire
flow works.

\subsubsection*{CSV Download Access Control}

The CSV download endpoint was tested by using an authenticated and an unauthenticated user.
For the unauthenticated part, GET request to \texttt{/api/download-csv} returned a 401
response and was blocked by the frontend, opening the login modal after. By logging in
correctly, the same request succeeded and the browser downloaded the complete sensor dataset
as a CSV file.

\subsubsection*{Data Page Filtering}

Multiple options of using the filter from the Data page have been tested. In all such cases
tested, the charts have been rendered with the correct filtered data set. The expected data
was filtered and shown. The Reset All button works as expected and removes all filters and
loads the complete, initial, dataset. If there is no data for the selected filters, then ``No
data found for the selected filters'' is shown.

\vspace{12pt}
\begin{figure}[H]
  \centering
  \includegraphics[width=0.75\textwidth]{data_filter}
  \caption{Figure 42: Data page filter}
  \label{fig:data-filter}
\end{figure}
\vspace{12pt}

\subsubsection*{Maintenance Mode}

Maintenance mode toggle was tested by enabling it from the Admin page. The website was
accessed afterwards by a non-admin session and only displayed the maintenance specific page.
Turning off maintenance mode, from the admin page, by the admin account user, restores the
normal access to the website for everyone. The administrator's own session is not affected.

\vspace{12pt}
\begin{figure}[H]
  \centering
  \includegraphics[width=0.75\textwidth]{maintenance_page}
  \caption{Figure 43: Maintenance page}
  \label{fig:maintenance}
\end{figure}
\vspace{12pt}

\subsubsection*{Android App}

It has been tested by deploying the APK build on a physical Android device to test its
connectivity with the Flask backend server. Home screen did show the images properly from the
asset files. Data screen fetched the sensor values and showed the line charts and the filter
dropdowns. Predictions screen showed the cards of the predictions with all the proper date
formatting, temperature range, humidity, weather indicator and danger part. The Contact screen
(form page on website) sent a test message to the \texttt{/api/contact} endpoint, which has
been found appearing in the Contact Messages table of the admin panel. The APK file download
was also tested from the Download page of the web application, the same way it is done for the
CSV file.

\section{Known Limitations}
\label{subsec:limitations}

There have been a number of issues identified during the testing that are documented for future
improvement.

If the room temperature is too high (30 or higher degrees Celsius), then the time necessary
for the sensor's initial heat up might need to be longer than the first 30 minutes before the
first reading after start.

The MQ6 gas sensor saturates the ADC in cases of high gas concentrations, compromising its
accuracy. In case of hardware modification, changing the value of the load resistor to reduce
the voltage range output will increase the concentration ceiling limit of measurement.

The formula used to correct the temperature and humidity effect on the PPM reading is a general
approximation, obtained from the Hanwei application note~\cite{hanwei}.

Currently, the setting for the neighbourhood can only be set at initialization through the
command terminal and once the station has been initialized, it cannot be changed without
stopping the whole process. The better approach will be having the setting stored in a
configuration file which is then read at initialization.

If the DHT11 sensor is placed too close to the other sensors, it will pick up their heat
instead. This will lead to wrong readings and the correction formula not being computed
correctly.

The MQ7 sensor is not very precise when it comes to very low readings, like the ones specific
to the actual carbon monoxide present in the indoor space. They give constant readings which
are 35\,ppm to 100\,ppm higher than they should, after heating. It is not something unusual
for this type of sensor.

The station doesn't have a portable electrical source available yet. It must be connected to
an electrical socket.

When it comes to the prediction system, even in case the system works perfectly fine,
predictions should be made with caution. It is also important to consider the length of the
dataset.
